\chapter{Body Dysmorphic Disorder}

\section*{Definition}
Body dysmorphic disorder (BDD) is characterized by a preoccupation with one or more perceived defects or flaws in physical appearance that are not observable or appear only slight to others. The preoccupation causes clinically significant distress or impairment and is not better explained by an eating disorder.

\section*{When to See a Doctor}

\begin{itemize}
 \item Appearance preoccupations cause significant distress or interfere with work, relationships, or daily life
 \item Repetitive behaviors (mirror checking, skin picking, camouflaging) are uncontrolled
 \item There are thoughts of suicide or self-harm; seek immediate emergency help and do not stay alone
\end{itemize}

\section*{How It Develops}
BDD has complex biological, psychological, and social contributors. Differences in visual processing, attention, emotion regulation, and threat perception have been reported, but no single brain abnormality or neurotransmitter explains the disorder.

\section*{Causes}
\begin{itemize}
 \item \textbf{Psychiatric:} Primary BDD (often classified with obsessive-compulsive and related disorders)
 \item \textbf{Genetic:} Family studies suggest heritability; elevated rates in first-degree relatives of BDD people
 \item \textbf{Environmental:} History of childhood teasing, abuse, or neglect; societal emphasis on appearance
 \item \textbf{Comorbidities:} Major depressive disorder, social anxiety disorder, OCD, substance use disorders, eating disorders
 \item \textbf{Psychological factors:} Perfectionism, aesthetic sensitivity, impaired emotion regulation
\end{itemize}

\section*{Related Symptoms}
\begin{itemize}
 \item Obsessive thoughts, compulsive checking, skin picking, camouflaging, social avoidance, depression, muscle dysmorphia
\end{itemize}
\textbf{Important:} Many individuals with BDD seek cosmetic or dermatologic treatment rather than psychiatric care, which rarely improves the preoccupation.

\section*{Medical Evaluation}
\begin{itemize}
 \item Clinical interview using DSM-5 criteria, supplemented by the BDD Diagnostic Module
 \item Severity assessment: Yale-Brown Obsessive-Compulsive Scale modified for BDD (BDD-YBOCS)
 \item Differentiate from eating disorders (weight/shape concerns), OCD (other obsessions), and gender dysphoria
 \item Screen for suicide risk, comorbid depression, and OCD
\end{itemize}

\section*{Treatment Options}
\begin{itemize}
 \item First-line psychotherapy is BDD-focused cognitive-behavioral therapy with exposure and response prevention.
 \item An SSRI, often fluoxetine, may be prescribed by a clinician; response can take several weeks and younger people or anyone at suicide risk needs close monitoring, especially after starting or changing the dose.
 \item If treatment does not help, specialist review may consider another SSRI, clomipramine, or carefully selected augmentation; antipsychotic combinations should not be started routinely in primary care.
 \item Cosmetic or dermatologic procedures usually do not treat the underlying preoccupation and should be considered only after appropriate assessment, while genuine medical concerns still deserve care.
\end{itemize}

\section*{Self-Care and Home Management}
\begin{itemize}
 \item Reduce mirror checking and reassurance-seeking behaviors gradually
 \item Practice exposure and response prevention with a qualified therapist rather than forcing difficult exercises alone.
 \item Keep a thought log to identify distorted beliefs about appearance
 \item Avoid comparing appearance to others, especially on social media
\end{itemize}

\section*{References}
\begin{thebibliography}{99}
\bibitem{body_dysmorphic_disorder:ref1} Phillips KA, Hollander E. ``Body dysmorphic disorder.'' \textit{N Engl J Med}. 2021;385(9):827-837.
\bibitem{body_dysmorphic_disorder:ref2} Veale D, Bewley A. ``Body dysmorphic disorder.'' \textit{BMJ}. 2015;350:h2278.
\bibitem{body_dysmorphic_disorder:ref3} Wilhelm S, Phillips KA, et al. ``Cognitive-behavioral therapy for body dysmorphic disorder.'' \textit{JAMA Psychiatry}. 2019;76(10):1077-1085.
\bibitem{body_dysmorphic_disorder:ref4} Sadock BJ, Sadock VA, et al. \textit{Kaplan \& Sadock\'s Comprehensive Textbook of Psychiatry}, 10th ed. Wolters Kluwer, 2017.
\bibitem{body_dysmorphic_disorder:ref5} National Institute for Health and Care Excellence. ``Obsessive-compulsive disorder and body dysmorphic disorder: treatment (CG31).'' Updated 2025. https://www.nice.org.uk/guidance/cg31.
\bibitem{body_dysmorphic_disorder:ref6} National Health Service. ``Body dysmorphic disorder (BDD).'' Accessed September 2026. https://www.nhs.uk/mental-health/conditions/body-dysmorphia/.
\end{thebibliography}
